\chapter{Decision Models}
\label{ch:models}

The system of \cref{ch:system} can execute any action; this chapter is about choosing which one.
Four answers are developed and compared: a rule-based reference policy, the \NS that is this
thesis's contribution, a deep reinforcement learning policy, and a family of ablations that
isolate what each part of the \NS actually contributes.

All of them implement the same interface. A policy receives the state of the factory and the
robot, together with the list of currently legal actions, and returns one action, a set of scores
and a human-readable explanation of why. Nothing in a policy touches ROS, which is what allows the
identical object to run inside the fast simulator during training and inside the decision service
on the robot (\cref{sec:service}).

\section{The rule-based reference}
\label{sec:rule-policy}

The reference policy exists so that the learned models have something honest to beat. It is
written from what an operator would say about the job, in five sentences:

\begin{enumerate}[leftmargin=*, itemsep=0.1em]
  \item Keep the robot able to get home: charge when the battery would not cover the trip plus the
        journey to the charger plus the reserve.
  \item Never let a line stand still: serve \textsc{blocked} sections first, then whichever will
        block soonest.
  \item Do not drive half empty: top the load up at another section if it is close and has units.
  \item Deliver what you carry when the robot is full or nothing else is worth collecting.
  \item Wait at the charger, never in an aisle, because waiting anywhere else only spends energy.
\end{enumerate}

A section counts as urgent if it will block within \SI{180}{\second}, and a buffer below
\SI{35}{\percent} is not worth a dedicated trip. The policy is roughly a hundred lines long, entirely
explainable, and --- as \cref{sec:overall} shows --- far from useless: it scores \num{41.66}
against the \NS's \num{53.81}, and it is the most energy-efficient policy in the comparison. It is
also the fallback the robot runs when the decision service is unavailable
(\cref{sec:service}), so the same code has to be trustworthy in production, not only in the
evaluation.

\section{The neuro-symbolic model}
\label{sec:ns-model}

The \NS is built from the observation that the decision has two different kinds of content. Some
of it is arithmetic and physics: whether the robot can complete a journey without stranding
itself, whether a line is already losing production, how much energy an action costs. That part is
known exactly, it can be computed from the robot's own models, and getting it wrong is
unacceptable rather than merely suboptimal. The rest is judgement: whether it is worth driving
past a half-full buffer to reach a fuller one, whether to collect a second load before delivering,
whether a quick top-up now saves a longer charge later. That part is not exactly computable, but
it can be learned.

The model separates the two. A symbolic layer decides what is \emph{allowed} and what is
\emph{urgent}; a neural scorer chooses among the candidates that survive. \Cref{fig:ns-pipeline}
shows the pipeline.

\begin{figure}[H]
    \centering
    \begin{tikzpicture}[
        font=\small,
        stage/.style={draw, rounded corners=2pt, align=center, minimum height=1.25cm,
                      text width=2.6cm, inner sep=4pt},
        symb/.style={stage, fill=black!6},
        neur/.style={stage, fill=black!12},
        io/.style={align=center, text width=2.3cm, inner sep=2pt},
        note/.style={font=\scriptsize, align=center, text width=2.8cm, text=black!65,
                     inner sep=2pt},
        flow/.style={-{Stealth[length=2mm]}, thick},
        drop/.style={-{Stealth[length=1.6mm]}, draw=black!55}]

        \node[io] (in) {legal\\actions};
        \node[symb, right=0.8cm of in]    (rules) {\textbf{hard rules}\\forbid};
        \node[symb, right=0.8cm of rules] (tiers) {\textbf{priority tiers}\\keep top tier};
        \node[neur, right=1.35cm of tiers] (rank)  {\textbf{ranking}\\estimate $+$\\bounded correction};
        \node[io, right=0.8cm of rank]    (out)   {action $+$\\explanation};

        \draw[flow] (in) -- (rules);
        \draw[flow] (rules) -- (tiers);
        \draw[flow] (tiers) -- (rank);
        \draw[flow] (rank) -- (out);

        % what each stage does, drawn downwards so the main flow stays a straight line
        \node[note, below=0.7cm of rules] (n1)
            {H1--H5: unsafe or\\impossible actions};
        \node[note, below=0.7cm of tiers] (n2)
            {tier 2 blocked line\\tier 1 about to block\\tier 0 routine};
        \node[note, below=0.7cm of rank]  (n3)
            {network may move an\\action by at most $\beta$};
        \draw[drop] (rules) -- (n1);
        \draw[drop] (tiers) -- (n2);
        \draw[drop] (rank)  -- (n3);

        \begin{scope}[on background layer]
            \node[draw, dashed, rounded corners, fit=(rules)(tiers)(n1)(n2),
                  inner sep=9pt,
                  label={[font=\scriptsize, text=black!65, inner sep=1pt, yshift=1pt]above:symbolic}] {};
            \node[draw, dashed, rounded corners, fit=(rank)(n3),
                  inner sep=9pt,
                  label={[font=\scriptsize, text=black!65, inner sep=1pt, yshift=1pt]above:neural}] {};
        \end{scope}
    \end{tikzpicture}
    \caption{The neuro-symbolic decision pipeline. The symbolic layer removes actions and ranks
             what remains into urgency tiers; the neural scorer only ever chooses \emph{within}
             the most urgent tier, and only by at most $\beta$ score units.}
    \label{fig:ns-pipeline}
\end{figure}

\subsection{Hard rules}
\label{sec:hard-rules}

Hard rules forbid. They encode the constraints of \cref{sec:constraints} and no score from any
network can overrule them --- this is the guarantee that a purely learned policy cannot give.

\begin{description}[leftmargin=1.2em, style=nextline, itemsep=0.2em]
  \item[H1 --- Maintenance.] Below \SI{30}{\percent} drive condition, pickups and deliveries are
        forbidden; only charging and waiting remain.
  \item[H2 --- Payload.] A pickup is forbidden if nothing at that section still fits within the
        remaining payload capacity.
  \item[H3 --- Reserve.] An action is forbidden if the battery after completing it \emph{and
        returning to the charger} would fall below the reserve plus a safety margin of
        \SI{2}{\percent}. This is not a threshold on the gauge: the whole route --- outward leg,
        dwell, journey home --- is simulated with the robot's own energy, thermal and wear models
        (\cref{sec:robot-model}), so the answer depends on the payload, the motor temperature and
        the condition of the drive.
  \item[H4 --- Heat.] An action is forbidden if the simulated peak motor temperature along that
        route would reach the limit of \SI{70}{\degreeCelsius}.
  \item[H5 --- Pointless waiting.] Away from the charger, if every pickup and delivery is already
        forbidden and charging is possible, waiting is forbidden too.
\end{description}

H1--H4 follow directly from the constraints. H5 does not: it was added after watching the robot
work a real shift, and the reason is worth recording. Charging becomes \emph{urgent} (tier~2
below) slightly below the battery level at which the pickups are already refused by H3. In the
gap between those two thresholds, waiting and charging looked equally good to the ranker, and the
robot spent about \SI{1}{\percent} of its pack waiting at a section before driving to the charger
anyway. The rule closes the gap by stating what is obviously true and was nowhere written down:
if the robot cannot collect or deliver anything until it charges, waiting cannot improve the
situation. Adding it improved the score on the training scenarios from \num{58.70} to \num{58.81}
and, on the full evaluation, \texttt{low\_battery\_start} from \num{55.14} to \num{55.57}.

\subsection{Priority tiers}
\label{sec:tiers}

The surviving actions are then ranked into three tiers, and only the most urgent tier is passed
on:

\begin{description}[leftmargin=1.2em, style=nextline, itemsep=0.2em]
  \item[Tier 2 --- unblock a stopped line.] A pickup at a \textsc{blocked} section, or charging
        when the battery is at the reserve.
  \item[Tier 1 --- prevent a line from stopping.] A pickup at a section that will fill within
        \SI{180}{\second} or within the travel time to it, whichever is longer; or a delivery when
        the robot is full and can collect nothing else.
  \item[Tier 0 --- routine.] Everything else.
\end{description}

The tiers express a priority that should never be traded away for a small gain in energy: a line
that is already losing production outranks one that is merely getting full, which outranks a
convenient errand. Because only the top tier reaches the scorer, the network can never choose a
routine action while a line stands still --- it can only decide \emph{which} of the equally urgent
actions to take.

\subsection{Explanation}
\label{sec:explanations}

Every verdict carries its reasons, so every decision explains itself. A real decision from a
Gazebo shift reads:

\begin{quote}\ttfamily\small
decision: WAIT:60 --- [routine] waiting away from the charger only spends energy |
PICKUP:B refused: battery would be 16 \% back at the charger, below the 17 \% reserve;
DELIVER refused: battery would be 17 \% back at the charger, below the 17 \% reserve;
neural score +0.95, chosen over CHARGE:40 by +0.16
\end{quote}

The explanation names the tier, the reason the chosen action is in it, which alternatives were
forbidden and why, and by how much the network preferred this action to the runner-up. This is not
a post-hoc rationalisation of a black-box output: it is the actual computation, written out. It is
also how the H5 rule was found --- the string above is the trace that showed the robot preferring
to wait while it was in fact stuck until it charged.

\section{The neural scorer}
\label{sec:neural-ranker}

\subsection{Features}
\label{sec:features}

The network scores \emph{one candidate action at a time}: it receives a vector of \num{27}
features describing that action in the current state, and returns a single number. Scoring actions
rather than states is what keeps the network small and portable --- it learns ``is this action
worth it, given what it costs and what it saves'', not ``what should be done in this exact
factory''.

The features cover four groups: what kind of action it is; the robot's state (battery,
temperature, condition, payload, cargo, time left in the shift); what the action costs and leaves
behind (distance, energy as a percentage of the pack, battery afterwards, battery after also
returning home); and what it achieves (the target section's fill, units and mass gained, time to
full, rate, health, whether it is blocked or faulted, plus the worst other section's fill and time
to full, and whether any other line is blocked). Charging actions additionally carry their target
level and the gain it would bring.

\subsection{Learning a correction, not a value}
\label{sec:residual}

The network does not predict an action's value. It predicts a \emph{correction} to an analytic
estimate:

\begin{equation}
    \label{eq:score}
    \operatorname{score}(a) \;=\; \underbrace{\hat{v}(a)}_{\text{symbolic estimate}}
        \;+\; \underbrace{c_\theta(a)}_{\text{learned correction}}
\end{equation}

where the symbolic estimate is the obvious part of the value, computed directly from the models:
the units the action moves, valued with the objective of \cref{eq:objective}, minus the energy the
journey and its dwell time cost. What the estimate cannot know is the future --- which line blocks
next, whether a fuller load would have been better, whether this trip makes a charge unavoidable
later --- and that is exactly what the correction has to supply.

This split matters for a reason that only became visible in the evaluation
(\cref{sec:bounded-correction}). The estimate is arithmetic on the current state, so it stays
correct in any factory, including one with twice the production rate. The network is a function
fitted on a distribution of states, so it is reliable near the situations it was trained on and
unreliable away from them. Starting from the estimate means the model degrades towards sensible
arithmetic rather than towards noise.

\subsection{Training}
\label{sec:ns-training}

Training data are collected by look-ahead in the fast simulator. At each decision point of a
simulated shift, every candidate action is taken in a copy of the simulator, the resulting future
is played out for \SI{900}{\second} with a reference policy, and the discounted return is recorded
as that action's value. Because production and faults are random, three futures are averaged per
action. Exploration is added by choosing a random action at \SI{25}{\percent} of decision points,
so the data are not confined to what the reference policy would have done. Collection runs over
the six training scenarios with seeds from \num{12345} upwards, which never overlap the evaluation
seeds (\cref{sec:splits}).

The decisive design choice is the loss. The look-ahead value of an action is noisy, and the
candidates at one decision point often differ by far less than that noise, so regressing values
spends the network's capacity on modelling randomness. What the model actually has to get right is
\emph{which candidate is best}. The scorer is therefore trained with a softmax cross-entropy over
the candidates of each decision point, with the look-ahead's best action as the target, plus a
small value-regression term that keeps the scores on a meaningful scale so they remain readable in
the explanation strings.

Training is repeated for three rounds of policy iteration: the first round looks ahead with the
symbolic policy, and each later round looks ahead with the model trained in the previous one. The
round with the best validation score is kept. In the final training run, round~1 was the best
(validation score \num{59.68}, against \num{58.05} and \num{56.10} for rounds~2 and~3), so policy
iteration did not pay off here --- a negative result worth reporting rather than hiding. The kept
model picks the look-ahead's best action in \SI{35}{\percent} of validation decision points, which
sounds low until one notices that many decision points have several near-equivalent candidates;
what matters is the shift score, not the agreement rate.

\section{Bounding the correction}
\label{sec:bounded-correction}

The model as described so far --- rules, tiers, and a network correcting an analytic estimate ---
is the version that was first evaluated. It was worse than its own symbolic layer on scenarios it
had never seen (\num{43.33} against \num{45.66}, \cref{tab:splits}), and the gap came almost
entirely from one scenario, \texttt{heavy\_load}, where it scores \num{\heavyUnbounded} against
the symbolic layer's \num{53.21}.\footnote{These are the numbers of the final configuration, in
which both variants also carry the H5 rule of \cref{sec:hard-rules}. When the failure was first
observed, before H5 existed, the same comparison read \num{\heavyDiscovery} against
\num{53.21}.}

The decision trace explained why. In a factory producing beyond one robot's capacity, the model
repeatedly chose \texttt{WAIT} while carrying cargo, up to eight times in a row, while all three
lines stood \textsc{blocked} and lost production. Its own explanation string said, correctly,
that ``waiting away from the charger only spends energy'': the symbolic estimate ranked waiting
last, and the learned correction --- around $+20$ on every candidate, with differences of a few
units --- simply outvoted it.

The cause is not a bug but a distribution shift. In the six training scenarios, production is
roughly half as fast, and waiting to accumulate a fuller load genuinely pays. The network learned
that habit and applied it to an overload it had never seen, where the same habit is ruinous.

The remedy follows from \cref{eq:score}. The correction is centred on the mean of the candidates
at that decision point --- only differences between candidates affect the ranking --- and then
clipped:

\begin{equation}
    \label{eq:bounded}
    \operatorname{score}(a) \;=\; \hat{v}(a)
        \;+\; \operatorname{clip}\!\left(c_\theta(a) - \overline{c_\theta},\; -\beta,\; +\beta\right)
\end{equation}

With a bound $\beta$, the network can reorder candidates whose analytic estimates lie within
$\beta$ of each other --- the close calls, where judgement is what is needed --- and can never
overturn a difference the arithmetic considers clear. It constrains the network's
\emph{influence}, which is a different thing from the action masking and shielding of
\cref{sec:safe-learning}: those constrain which actions are \emph{legal}.

The bound was selected on the six training scenarios with validation seeds \numrange{1000}{1019},
never on the unseen scenarios or the evaluation seeds (\cref{tab:bound-sweep}).

\begin{table}[H]
    \centering
    \caption{Choosing the bound $\beta$: mean score over the six training scenarios with
             validation seeds \numrange{1000}{1019}. $\beta=\infty$ is the unbounded model,
             $\beta=0$ switches the network off and decides on the analytic estimate alone.}
    \label{tab:bound-sweep}
    \begin{tabular}{lrrrrrrr}
        \toprule
        $\beta$ & $\infty$ & 20 & 10 & 5 & 2 & \textbf{1} & 0\\
        \midrule
        Mean score & 57.75 & 57.75 & 58.06 & 58.15 & 58.15 & \textbf{58.70} & 57.41\\
        \bottomrule
    \end{tabular}
\end{table}

Two things are visible in that table. The network is worth having: switching it off entirely
($\beta=0$) is worse than every bounded setting. And its useful influence is small: the best value
is $\beta=\maxCorrection$, one point on a scale where a delivered unit is worth one point. The
chosen model may therefore move an action by at most the value of a single delivered unit relative
to the others --- enough to break a tie, not enough to argue with the arithmetic.

What this costs and buys is measured in \cref{sec:ablation-results}: nothing on the training
scenarios ($p=0.058$), \num{3.87} points on the unseen ones ($p<0.001$), and
\num{\heavyBounded} against \num{\heavyUnbounded} on \texttt{heavy\_load}.

\section{The reinforcement learning comparison}
\label{sec:ppo}

The learned comparison is Maskable PPO, trained on the same environment, the same observation
content and the same objective, so that the comparison is about \emph{where knowledge is placed},
not about who gets more information.

The agent sees the state as one vector (robot state plus per-section features) and chooses from
the same discrete action set. Action masks mark the actions that are currently legal, which is the
standard way of expressing hard constraints in reinforcement learning, and the agent samples only
among them. Training runs for \num{300000} steps across four parallel environments
($n_{\text{steps}}=512$, batch size \num{256}, learning rate \num{3e-4}, $\gamma=0.995$,
$\lambda=0.95$, entropy coefficient \num{0.01}, two hidden layers of \num{64} units --- the same
size as the neural scorer), and takes about ten minutes on the evaluation machine.

One detail of the training setup had to be corrected before the comparison was fair. Each parallel
environment was originally bound to one fixed scenario, so with four environments the agent only
ever saw four of the six training scenarios. The environment now draws a scenario at random at
every reset, so all six are seen regardless of how many environments run. Both the corrected model
and the original one are reported (\texttt{ppo} and \texttt{ppo\_wp6}), because the difference
between them is itself informative: it is small, and the safety behaviour does not change
(\cref{sec:safety-results}).

A second variant, \texttt{ppo\_unmasked}, is trained without the action mask. Illegal choices are
executed anyway --- they simply achieve nothing --- and cost a small penalty, so the agent has to
learn the rules from experience. It answers the question of how much of a masked agent's
competence comes from the mask rather than from learning.

\section{Ablations}
\label{sec:ablations}

\Cref{tab:policies} lists every policy in the evaluation and the question it exists to answer.
Each removes exactly one thing, so that a difference in \cref{ch:results} can be attributed.

\begin{table}[H]
    \centering
    \caption{The policies compared in \cref{ch:results} and what each one isolates.}
    \label{tab:policies}
    \begin{tabularx}{\textwidth}{lX}
        \toprule
        Policy & What it is, and the question it answers\\
        \midrule
        \texttt{ns} & The full model: hard rules, priority tiers, analytic estimate, bounded neural correction.\\
        \texttt{ns\_unbounded} & The same model with $\beta=\infty$. \emph{Does bounding the network matter?}\\
        \texttt{ns\_estimate\_only} & The same model with $\beta=0$. \emph{Does the network contribute anything at all?}\\
        \texttt{ns\_symbolic\_only} & Rules and tiers, then the cheapest action per unit moved; no network, no estimate. \emph{How far do the rules alone get?}\\
        \texttt{ns\_neural\_only} & The network ranking every legal action, with no rules and no tiers. \emph{What did the network really learn --- including about safety?}\\
        \texttt{rule} & The reference policy of \cref{sec:rule-policy}. \emph{Is any of this better than common sense?}\\
        \texttt{ppo} & Maskable PPO trained on all six training scenarios. \emph{Rules versus learning, with equal information.}\\
        \texttt{ppo\_unmasked} & PPO without the action mask. \emph{How much of PPO's competence is the mask?}\\
        \texttt{ppo\_wp6} & The earlier PPO, trained on four of the six scenarios. \emph{Did the training-coverage fix change the comparison?}\\
        \bottomrule
    \end{tabularx}
\end{table}
